% documentation_EN.tex — WAC (Windows Artefacts Collector) user manual, English edition.
% The French edition is documentation_FR.tex, in this folder: both are maintained
% together, section for section.
%
% Build (XeLaTeX, two passes for the cover page and the table of contents):
%     latexmk -xelatex documentation_EN.tex
% The cover page is ../cover_EN.png, generated from the French one by
% ../translate-cover.py.

\documentclass[11pt,a4paper,oneside]{report}

% ---------------------------------------------------------------- language and labels
\usepackage[british]{babel}
\newcommand{\wacdoctitle}{User manual}
\newcommand{\wacpageof}{of}
\newcommand{\wacnotetitle}{Keep in mind}
\newcommand{\wacwarningtitle}{Warning}
% Layout shared with every WAC document.
% wac-style.tex — layout shared by every WAC document (user manuals, testing
% guides), French and English editions alike: one copy, so that they cannot
% drift apart.
%
% A document sets, BEFORE % wac-style.tex — layout shared by every WAC document (user manuals, testing
% guides), French and English editions alike: one copy, so that they cannot
% drift apart.
%
% A document sets, BEFORE % wac-style.tex — layout shared by every WAC document (user manuals, testing
% guides), French and English editions alike: one copy, so that they cannot
% drift apart.
%
% A document sets, BEFORE \input{../wac-style.tex}:
%   \usepackage[<language>]{babel}
%   \newcommand{\wacdoctitle}{...}      running header: "WAC — <title>"
%   \newcommand{\wacpageof}{...}        footer: "Page N <of> M"
%   \newcommand{\wacnotetitle}{...}     default title of the "note" box
%   \newcommand{\wacwarningtitle}{...}  default title of the "attention" box
% and AFTER it its \hypersetup (PDF title, author, subject).

% ---------------------------------------------------------------- fonts, language
\usepackage{fontspec}
\setmainfont{Noto Sans}
\setmonofont{DejaVu Sans Mono}[Scale=0.84]

% ---------------------------------------------------------------- layout
\usepackage[top=26mm,bottom=26mm,left=24mm,right=22mm,headheight=14pt]{geometry}
\usepackage{graphicx}
\graphicspath{{../}}
\usepackage{xcolor}
\usepackage{tikz}
\usepackage{booktabs,tabularx,longtable,array}
\usepackage{enumitem}
\usepackage{fancyhdr}
\usepackage{titlesec}
\usepackage[most]{tcolorbox}
\usepackage{listings}
\usepackage{url}
\usepackage{lastpage}
\usepackage[hidelinks,unicode,pdfpagelabels]{hyperref}

% Colours sampled from the cover page
\definecolor{marine}{HTML}{0A1A33}
\definecolor{bleu}{HTML}{2F6DB5}
\definecolor{accent}{HTML}{1F66D1}
\definecolor{fond}{HTML}{EEF3FA}
\definecolor{alerte}{HTML}{B3261E}
\definecolor{fondalerte}{HTML}{FBEDEC}
\definecolor{gris}{HTML}{5B6573}

\setlength{\parindent}{0pt}
\setlength{\parskip}{0.55em}
\linespread{1.08}
\setlist{itemsep=0.25em,topsep=0.35em}
\urlstyle{tt}
\renewcommand{\arraystretch}{1.25}
% Left-aligned table columns: justified narrow columns opened wide gaps
% between words.
\newcolumntype{L}[1]{>{\raggedright\arraybackslash}p{#1}}
% Monospaced file and field names cannot be hyphenated: let the line's spaces
% stretch instead.
\setlength{\emergencystretch}{3em}

% Headings
\titleformat{\chapter}[display]
  {\color{marine}\bfseries\raggedright}
  {\color{bleu}\large\MakeUppercase{\chaptertitlename}~\thechapter}
  {0.3em}
  {\Huge}[{\color{accent}\titlerule[2pt]}]
\titlespacing*{\chapter}{0pt}{0pt}{1.6em}
\titleformat{\section}{\color{marine}\Large\bfseries}{\color{bleu}\thesection}{0.8em}{}
\titleformat{\subsection}{\color{marine}\large\bfseries}{\color{bleu}\thesubsection}{0.7em}{}
\titleformat{\subsubsection}{\color{bleu}\normalsize\bfseries}{}{0em}{}

% Headers and footers
\pagestyle{fancy}
\fancyhf{}
\fancyhead[L]{\small\color{gris}WAC — \wacdoctitle}
\fancyhead[R]{\small\color{gris}\leftmark}
% Footer: "Page N of M" (the word set by \wacpageof), readable. A bare, small, light-grey number went
% unnoticed.
\newcommand{\piedpage}{\small\color{marine}\textbf{Page \thepage}\ \wacpageof\ \pageref*{LastPage}}
\fancyfoot[C]{\piedpage}
\renewcommand{\footrulewidth}{0.4pt}
\renewcommand{\footrule}{\hbox to\headwidth{\color{fond}\leaders\hrule height \footrulewidth\hfill}}
\renewcommand{\headrulewidth}{0.4pt}
\renewcommand{\headrule}{\hbox to\headwidth{\color{fond}\leaders\hrule height \headrulewidth\hfill}}
\renewcommand{\chaptermark}[1]{\markboth{#1}{}}
\fancypagestyle{plain}{\fancyhf{}\fancyfoot[C]{\piedpage}\renewcommand{\headrulewidth}{0pt}}

% Boxes
\newtcolorbox{note}[1][\wacnotetitle]{enhanced,breakable,colback=fond,colframe=bleu,
  boxrule=0pt,leftrule=3pt,arc=0pt,left=8pt,right=8pt,top=6pt,bottom=6pt,
  title={\bfseries #1},coltitle=bleu,colbacktitle=fond,attach title to upper={\par\smallskip}}
\newtcolorbox{attention}[1][\wacwarningtitle]{enhanced,breakable,colback=fondalerte,colframe=alerte,
  boxrule=0pt,leftrule=3pt,arc=0pt,left=8pt,right=8pt,top=6pt,bottom=6pt,
  title={\bfseries #1},coltitle=alerte,colbacktitle=fondalerte,attach title to upper={\par\smallskip}}

% Code
\lstset{basicstyle=\ttfamily\small,columns=fullflexible,keepspaces=true,
  backgroundcolor=\color{fond},frame=l,framerule=3pt,rulecolor=\color{bleu},
  xleftmargin=9pt,framexleftmargin=6pt,aboveskip=0.8em,belowskip=0.8em,
  breaklines=true,breakatwhitespace=false,showstringspaces=false,
  literate={é}{{\'e}}1 {è}{{\`e}}1 {à}{{\`a}}1 {ê}{{\^e}}1 {ô}{{\^o}}1 {ç}{{\c c}}1
           {É}{{\'E}}1 {—}{{---}}1 {’}{{'}}1 {…}{{...}}1 {«}{{\guillemotleft}}1 {»}{{\guillemotright}}1}

% Shortcuts
\newcommand{\wac}{\textsc{wac}}
\newcommand{\opt}[1]{\texttt{-{}-#1}}
\newcommand{\champ}[1]{\texttt{#1}}
\newcommand{\fichier}[1]{\texttt{#1}}

% ============================================================ COVER PAGE
% The cover images are 2:3 (1024 x 1536), taller than an A4 page once widened
% to it: centred, they lost 9 mm at the top — the CEC logo — and 9 mm at the
% bottom. Raised 4 mm above the top of the page instead: the logo, which starts
% 6 mm below the image's edge, stays whole, and the bottom row of text keeps
% 7 mm to the page's edge; only a decorative corner below it is cut.
% Usage: \coverpage{cover_FR.png}
\newcommand{\coverpage}[1]{%
  \begin{titlepage}
  \thispagestyle{empty}
  \begin{tikzpicture}[remember picture,overlay]
    \node[inner sep=0pt,anchor=north,yshift=4mm] at (current page.north)
      {\includegraphics[width=\paperwidth]{#1}};
  \end{tikzpicture}
  \end{titlepage}}
:
%   \usepackage[<language>]{babel}
%   \newcommand{\wacdoctitle}{...}      running header: "WAC — <title>"
%   \newcommand{\wacpageof}{...}        footer: "Page N <of> M"
%   \newcommand{\wacnotetitle}{...}     default title of the "note" box
%   \newcommand{\wacwarningtitle}{...}  default title of the "attention" box
% and AFTER it its \hypersetup (PDF title, author, subject).

% ---------------------------------------------------------------- fonts, language
\usepackage{fontspec}
\setmainfont{Noto Sans}
\setmonofont{DejaVu Sans Mono}[Scale=0.84]

% ---------------------------------------------------------------- layout
\usepackage[top=26mm,bottom=26mm,left=24mm,right=22mm,headheight=14pt]{geometry}
\usepackage{graphicx}
\graphicspath{{../}}
\usepackage{xcolor}
\usepackage{tikz}
\usepackage{booktabs,tabularx,longtable,array}
\usepackage{enumitem}
\usepackage{fancyhdr}
\usepackage{titlesec}
\usepackage[most]{tcolorbox}
\usepackage{listings}
\usepackage{url}
\usepackage{lastpage}
\usepackage[hidelinks,unicode,pdfpagelabels]{hyperref}

% Colours sampled from the cover page
\definecolor{marine}{HTML}{0A1A33}
\definecolor{bleu}{HTML}{2F6DB5}
\definecolor{accent}{HTML}{1F66D1}
\definecolor{fond}{HTML}{EEF3FA}
\definecolor{alerte}{HTML}{B3261E}
\definecolor{fondalerte}{HTML}{FBEDEC}
\definecolor{gris}{HTML}{5B6573}

\setlength{\parindent}{0pt}
\setlength{\parskip}{0.55em}
\linespread{1.08}
\setlist{itemsep=0.25em,topsep=0.35em}
\urlstyle{tt}
\renewcommand{\arraystretch}{1.25}
% Left-aligned table columns: justified narrow columns opened wide gaps
% between words.
\newcolumntype{L}[1]{>{\raggedright\arraybackslash}p{#1}}
% Monospaced file and field names cannot be hyphenated: let the line's spaces
% stretch instead.
\setlength{\emergencystretch}{3em}

% Headings
\titleformat{\chapter}[display]
  {\color{marine}\bfseries\raggedright}
  {\color{bleu}\large\MakeUppercase{\chaptertitlename}~\thechapter}
  {0.3em}
  {\Huge}[{\color{accent}\titlerule[2pt]}]
\titlespacing*{\chapter}{0pt}{0pt}{1.6em}
\titleformat{\section}{\color{marine}\Large\bfseries}{\color{bleu}\thesection}{0.8em}{}
\titleformat{\subsection}{\color{marine}\large\bfseries}{\color{bleu}\thesubsection}{0.7em}{}
\titleformat{\subsubsection}{\color{bleu}\normalsize\bfseries}{}{0em}{}

% Headers and footers
\pagestyle{fancy}
\fancyhf{}
\fancyhead[L]{\small\color{gris}WAC — \wacdoctitle}
\fancyhead[R]{\small\color{gris}\leftmark}
% Footer: "Page N of M" (the word set by \wacpageof), readable. A bare, small, light-grey number went
% unnoticed.
\newcommand{\piedpage}{\small\color{marine}\textbf{Page \thepage}\ \wacpageof\ \pageref*{LastPage}}
\fancyfoot[C]{\piedpage}
\renewcommand{\footrulewidth}{0.4pt}
\renewcommand{\footrule}{\hbox to\headwidth{\color{fond}\leaders\hrule height \footrulewidth\hfill}}
\renewcommand{\headrulewidth}{0.4pt}
\renewcommand{\headrule}{\hbox to\headwidth{\color{fond}\leaders\hrule height \headrulewidth\hfill}}
\renewcommand{\chaptermark}[1]{\markboth{#1}{}}
\fancypagestyle{plain}{\fancyhf{}\fancyfoot[C]{\piedpage}\renewcommand{\headrulewidth}{0pt}}

% Boxes
\newtcolorbox{note}[1][\wacnotetitle]{enhanced,breakable,colback=fond,colframe=bleu,
  boxrule=0pt,leftrule=3pt,arc=0pt,left=8pt,right=8pt,top=6pt,bottom=6pt,
  title={\bfseries #1},coltitle=bleu,colbacktitle=fond,attach title to upper={\par\smallskip}}
\newtcolorbox{attention}[1][\wacwarningtitle]{enhanced,breakable,colback=fondalerte,colframe=alerte,
  boxrule=0pt,leftrule=3pt,arc=0pt,left=8pt,right=8pt,top=6pt,bottom=6pt,
  title={\bfseries #1},coltitle=alerte,colbacktitle=fondalerte,attach title to upper={\par\smallskip}}

% Code
\lstset{basicstyle=\ttfamily\small,columns=fullflexible,keepspaces=true,
  backgroundcolor=\color{fond},frame=l,framerule=3pt,rulecolor=\color{bleu},
  xleftmargin=9pt,framexleftmargin=6pt,aboveskip=0.8em,belowskip=0.8em,
  breaklines=true,breakatwhitespace=false,showstringspaces=false,
  literate={é}{{\'e}}1 {è}{{\`e}}1 {à}{{\`a}}1 {ê}{{\^e}}1 {ô}{{\^o}}1 {ç}{{\c c}}1
           {É}{{\'E}}1 {—}{{---}}1 {’}{{'}}1 {…}{{...}}1 {«}{{\guillemotleft}}1 {»}{{\guillemotright}}1}

% Shortcuts
\newcommand{\wac}{\textsc{wac}}
\newcommand{\opt}[1]{\texttt{-{}-#1}}
\newcommand{\champ}[1]{\texttt{#1}}
\newcommand{\fichier}[1]{\texttt{#1}}

% ============================================================ COVER PAGE
% The cover images are 2:3 (1024 x 1536), taller than an A4 page once widened
% to it: centred, they lost 9 mm at the top — the CEC logo — and 9 mm at the
% bottom. Raised 4 mm above the top of the page instead: the logo, which starts
% 6 mm below the image's edge, stays whole, and the bottom row of text keeps
% 7 mm to the page's edge; only a decorative corner below it is cut.
% Usage: \coverpage{cover_FR.png}
\newcommand{\coverpage}[1]{%
  \begin{titlepage}
  \thispagestyle{empty}
  \begin{tikzpicture}[remember picture,overlay]
    \node[inner sep=0pt,anchor=north,yshift=4mm] at (current page.north)
      {\includegraphics[width=\paperwidth]{#1}};
  \end{tikzpicture}
  \end{titlepage}}
:
%   \usepackage[<language>]{babel}
%   \newcommand{\wacdoctitle}{...}      running header: "WAC — <title>"
%   \newcommand{\wacpageof}{...}        footer: "Page N <of> M"
%   \newcommand{\wacnotetitle}{...}     default title of the "note" box
%   \newcommand{\wacwarningtitle}{...}  default title of the "attention" box
% and AFTER it its \hypersetup (PDF title, author, subject).

% ---------------------------------------------------------------- fonts, language
\usepackage{fontspec}
\setmainfont{Noto Sans}
\setmonofont{DejaVu Sans Mono}[Scale=0.84]

% ---------------------------------------------------------------- layout
\usepackage[top=26mm,bottom=26mm,left=24mm,right=22mm,headheight=14pt]{geometry}
\usepackage{graphicx}
\graphicspath{{../}}
\usepackage{xcolor}
\usepackage{tikz}
\usepackage{booktabs,tabularx,longtable,array}
\usepackage{enumitem}
\usepackage{fancyhdr}
\usepackage{titlesec}
\usepackage[most]{tcolorbox}
\usepackage{listings}
\usepackage{url}
\usepackage{lastpage}
\usepackage[hidelinks,unicode,pdfpagelabels]{hyperref}

% Colours sampled from the cover page
\definecolor{marine}{HTML}{0A1A33}
\definecolor{bleu}{HTML}{2F6DB5}
\definecolor{accent}{HTML}{1F66D1}
\definecolor{fond}{HTML}{EEF3FA}
\definecolor{alerte}{HTML}{B3261E}
\definecolor{fondalerte}{HTML}{FBEDEC}
\definecolor{gris}{HTML}{5B6573}

\setlength{\parindent}{0pt}
\setlength{\parskip}{0.55em}
\linespread{1.08}
\setlist{itemsep=0.25em,topsep=0.35em}
\urlstyle{tt}
\renewcommand{\arraystretch}{1.25}
% Left-aligned table columns: justified narrow columns opened wide gaps
% between words.
\newcolumntype{L}[1]{>{\raggedright\arraybackslash}p{#1}}
% Monospaced file and field names cannot be hyphenated: let the line's spaces
% stretch instead.
\setlength{\emergencystretch}{3em}

% Headings
\titleformat{\chapter}[display]
  {\color{marine}\bfseries\raggedright}
  {\color{bleu}\large\MakeUppercase{\chaptertitlename}~\thechapter}
  {0.3em}
  {\Huge}[{\color{accent}\titlerule[2pt]}]
\titlespacing*{\chapter}{0pt}{0pt}{1.6em}
\titleformat{\section}{\color{marine}\Large\bfseries}{\color{bleu}\thesection}{0.8em}{}
\titleformat{\subsection}{\color{marine}\large\bfseries}{\color{bleu}\thesubsection}{0.7em}{}
\titleformat{\subsubsection}{\color{bleu}\normalsize\bfseries}{}{0em}{}

% Headers and footers
\pagestyle{fancy}
\fancyhf{}
\fancyhead[L]{\small\color{gris}WAC — \wacdoctitle}
\fancyhead[R]{\small\color{gris}\leftmark}
% Footer: "Page N of M" (the word set by \wacpageof), readable. A bare, small, light-grey number went
% unnoticed.
\newcommand{\piedpage}{\small\color{marine}\textbf{Page \thepage}\ \wacpageof\ \pageref*{LastPage}}
\fancyfoot[C]{\piedpage}
\renewcommand{\footrulewidth}{0.4pt}
\renewcommand{\footrule}{\hbox to\headwidth{\color{fond}\leaders\hrule height \footrulewidth\hfill}}
\renewcommand{\headrulewidth}{0.4pt}
\renewcommand{\headrule}{\hbox to\headwidth{\color{fond}\leaders\hrule height \headrulewidth\hfill}}
\renewcommand{\chaptermark}[1]{\markboth{#1}{}}
\fancypagestyle{plain}{\fancyhf{}\fancyfoot[C]{\piedpage}\renewcommand{\headrulewidth}{0pt}}

% Boxes
\newtcolorbox{note}[1][\wacnotetitle]{enhanced,breakable,colback=fond,colframe=bleu,
  boxrule=0pt,leftrule=3pt,arc=0pt,left=8pt,right=8pt,top=6pt,bottom=6pt,
  title={\bfseries #1},coltitle=bleu,colbacktitle=fond,attach title to upper={\par\smallskip}}
\newtcolorbox{attention}[1][\wacwarningtitle]{enhanced,breakable,colback=fondalerte,colframe=alerte,
  boxrule=0pt,leftrule=3pt,arc=0pt,left=8pt,right=8pt,top=6pt,bottom=6pt,
  title={\bfseries #1},coltitle=alerte,colbacktitle=fondalerte,attach title to upper={\par\smallskip}}

% Code
\lstset{basicstyle=\ttfamily\small,columns=fullflexible,keepspaces=true,
  backgroundcolor=\color{fond},frame=l,framerule=3pt,rulecolor=\color{bleu},
  xleftmargin=9pt,framexleftmargin=6pt,aboveskip=0.8em,belowskip=0.8em,
  breaklines=true,breakatwhitespace=false,showstringspaces=false,
  literate={é}{{\'e}}1 {è}{{\`e}}1 {à}{{\`a}}1 {ê}{{\^e}}1 {ô}{{\^o}}1 {ç}{{\c c}}1
           {É}{{\'E}}1 {—}{{---}}1 {’}{{'}}1 {…}{{...}}1 {«}{{\guillemotleft}}1 {»}{{\guillemotright}}1}

% Shortcuts
\newcommand{\wac}{\textsc{wac}}
\newcommand{\opt}[1]{\texttt{-{}-#1}}
\newcommand{\champ}[1]{\texttt{#1}}
\newcommand{\fichier}[1]{\texttt{#1}}

% ============================================================ COVER PAGE
% The cover images are 2:3 (1024 x 1536), taller than an A4 page once widened
% to it: centred, they lost 9 mm at the top — the CEC logo — and 9 mm at the
% bottom. Raised 4 mm above the top of the page instead: the logo, which starts
% 6 mm below the image's edge, stays whole, and the bottom row of text keeps
% 7 mm to the page's edge; only a decorative corner below it is cut.
% Usage: \coverpage{cover_FR.png}
\newcommand{\coverpage}[1]{%
  \begin{titlepage}
  \thispagestyle{empty}
  \begin{tikzpicture}[remember picture,overlay]
    \node[inner sep=0pt,anchor=north,yshift=4mm] at (current page.north)
      {\includegraphics[width=\paperwidth]{#1}};
  \end{tikzpicture}
  \end{titlepage}}


\hypersetup{pdftitle={WAC — User manual},
  pdfauthor={Aerospace Cyber Defence Centre of Excellence (CEC)},
  pdfsubject={Windows artefact collection software for digital forensic investigation}}

\begin{document}

% ============================================================ COVER PAGE
\pagenumbering{roman}
\coverpage{cover_EN.png}

% ============================================================ INFORMATION PAGE
\thispagestyle{empty}
\vspace*{\fill}
{\color{marine}\Large\bfseries WAC — Windows Artefacts Collector}\par
{\color{bleu}User manual}\par
\medskip
\begin{tabular}{@{}ll@{}}
Software version & 1.4.0 \\
Document edition & September 2026 \\
Publisher & Aerospace Cyber Defence Centre of Excellence (CEC) \\
          & French Air and Space Force Academy, Air Base 701, Salon-de-Provence \\
\end{tabular}

\bigskip
{\small\color{gris}This document describes how to use \wac{}: preparing and
running a collection, reading and verifying its results, and the traces it
leaves on the examined machine. The figures quoted were measured on a
Windows~11 test virtual machine; they are orders of magnitude, not
commitments.}
\vspace*{2cm}
\clearpage

% Numbering starts at 1 with the table of contents; only the cover page and the
% information page, which have no footer, are left out of it.
\pagenumbering{arabic}
\tableofcontents

% ============================================================ 1
\chapter{Overview}

\section{What WAC is for}

\wac{} collects, from a running Windows computer, the artefacts a digital
investigation needs: registry, event logs, Prefetch files, shortcuts and jump
lists, scheduled tasks, services, accounts, processes and sessions. It delivers
them as \textbf{JSON} files, ready to be analysed or fed into a processing
pipeline.

It runs from a USB stick, without installation, and writes everything it
produces to that stick.

\section{The guiding principle: a minimal footprint}

Examining a machine changes it. Every tool leaves traces: opened files whose
last-access date changes, services that write to the logs when solicited,
registry keys read through the API. These traces mix with the ones being looked
for.

\wac{} is designed so that what it leaves behind is as close to nothing as the
task allows, and so that \textbf{whatever it cannot avoid is written down}.
In practice:

\begin{itemize}
  \item the disk is read \textbf{directly}, read-only, without the system
        opening any file: no access date is modified;
  \item no key of the machine's registry is opened: the hives are copied, then
        read on the copy;
  \item no Volume Shadow Copy, no COM or WMI call, no logging service is
        solicited;
  \item nothing is written to the examined disk;
  \item every operation is recorded, with the trace it is expected to leave, in
        \fichier{investigation.json}.
\end{itemize}

\section{What WAC does not do}

\begin{itemize}
  \item \wac{} \textbf{does not analyse}: it collects and formats. Interpretation
        is the analyst's.
  \item It does not make a full disk image: it extracts the files that hold
        artefacts.
  \item It does not capture memory. The only volatile data collected are the
        process list, the open sessions and the state of the services.
\end{itemize}

% ============================================================ 2
\chapter{How it works}

\section{Reading the disk directly}

\wac{} opens the volume (\path|\\.\C:|) read-only, walks the NTFS file table
(\path|$MFT|) itself and extracts the content of the files it needs. The system
sees no file being opened: no access date updated, no directory walked, no
lock.

This reading handles the NTFS features Windows normally hides:

\begin{description}[style=nextline,leftmargin=1.2em]
  \item[NTFS compression] Windows~11 compresses, among others, the event log
        folder. \wac{} decompresses it itself.
  \item[Compact OS (WOF)] The system binaries of Windows~10 and~11 are stored
        compressed in a hidden stream; seen through the API, they look ordinary.
        \wac{} decompresses them (XPRESS variants).
  \item[Split attributes] A heavily fragmented file has its information spread
        over several records; \wac{} gathers them.
  \item[Preallocated files] Beyond the valid data length, NTFS returns zeros,
        even if the disk still holds remnants of former files. \wac{} does the
        same: the extracted content is the one any ordinary acquisition would
        see, and so is its fingerprint.
\end{description}

\section{Exhibit store and working copy}

A digital exhibit is never examined on itself: a copy is taken and sealed, and
the work is done on a second copy. If the analysis damages something, the sealed
copy remains and the operation can be redone.

\wac{} applies this procedure to \textbf{every} extracted file, including those
it does not modify. The two directories keep their French names in the output:

\begin{description}[style=nextline,leftmargin=1.2em]
  \item[\fichier{exhibits/} — the exhibit store] the raw copies, as read from
        the volume. This directory is never reopened for writing after
        extraction. It holds the manifest identifying each exhibit, and the
        manifest's seal.
  \item[\fichier{working/} — the working copy] copied from the exhibit store and
        verified by fingerprint. The only modifications (replay of the hives'
        transaction logs) are made here, and every collector reads from here.
\end{description}

\begin{note}
Every extracted file is therefore written \textbf{twice} to the collection
medium. Plan for the space (see section~\ref{sec:place}).
\end{note}

\section{Registry hives and transaction logs}

A hive copied from a running machine lacks its latest changes: they are still in
its transaction logs (\fichier{.LOG1}, \fichier{.LOG2}). \wac{} extracts these
logs and \textbf{replays them onto the working copy}, so that the analysis bears
on the machine's actual state.

Everything \wac{} writes onto a copy can be undone: before the replay, the
original content of every replaced page is saved in an undo journal
(\fichier{<hive>.undo}). The raw copy stays intact in the exhibit store.

\section{Two-pass extraction}

Where each user's hives live is recorded in the \fichier{SOFTWARE} hive. \wac{}
therefore first extracts the machine's hives (\fichier{SYSTEM},
\fichier{SOFTWARE}, \fichier{SAM}, \fichier{Amcache}), reads the profile list in
the copy of \fichier{SOFTWARE}, then extracts each profile's hives
(\fichier{ntuser.dat}, \fichier{usrClass.dat}).

Nothing is assumed about \fichier{C:}: the system drive is detected, and a
profile on another disk (for instance \path|D:\Users\...|) is read on its own
volume.

\section{What is still read live}

A few pieces of information describe the very moment of the collection and
are held in no file. They are read live, with a single query each, and
recorded as read in the exhibit store (\fichier{exhibits\textbackslash live}),
where the conversion reads them back:

\begin{itemize}
  \item the running \textbf{processes} (a kernel snapshot);
  \item the open \textbf{sessions};
  \item the \textbf{current state of the services} (their configuration is read
        from the copied hive);
  \item the \textbf{clock} and the boot time;
  \item the running machine's \textbf{settings}: system drive, code page, time
        zone, volumes and their serial numbers.
\end{itemize}

% ============================================================ 3
\chapter{Preparing a collection}

\section{Requirements}

\begin{itemize}
  \item Windows~10 or later, 64-bit. \wac{} checks it at start-up.
  \item \textbf{Administrator rights}: reading a volume directly requires them.
        \wac{} checks this too.
  \item \textbf{NTFS}, \textbf{FAT} or \textbf{exFAT} volumes. A file on another
        file system (CD-ROM) cannot be read directly; it is reported as unread,
        never opened another way.
\end{itemize}

\section{The collection medium}\label{sec:place}

\wac{} runs from the collection medium (usually a USB stick) and writes all its
results there. Choose an \textbf{empty} medium, preferably USB~3: writing is the
most expensive part of a collection.

The space needed depends on the mode and the options. Two places fill up: the
\textbf{exhibit store} (sealed raw copies) and the \textbf{working copy}, which
takes only what the collectors modify or read during the collection (hives,
Prefetch, shortcuts, jump lists, referenced binaries). The binaries acquired by
\opt{collect} \opt{binary} and the NTFS metadata are never copied again: the
conversion reads them in the exhibit store itself. Measured on the test virtual
machine (Windows~11 without added software, 1~volume):

\begin{center}
\begin{tabular}{@{}L{7.4cm}rrr@{}}
\toprule
Mode & Exhibit store & Working copy & Time \\
\midrule
Full run, \opt{events} \opt{binary} (referenced binaries) & 1.2~GB & 0.6~GB & 3~min~30 \\
Collection only, \opt{collect} \opt{events} \opt{binary} (every executable) & 2.7~GB & 0.3~GB & 8~min~30 \\
\bottomrule
\end{tabular}
\end{center}

The heaviest items of the exhibit store:

\begin{center}
\begin{tabular}{@{}L{9.6cm}r@{}}
\toprule
Content & Exhibit store \\
\midrule
NTFS metadata: \path|$MFT| (380~MB), \path|$LogFile| (64~MB),
  \path|$UsnJrnl:$J| (35~MB, out of 3.6~GB logical: NTFS freed the start,
  which is not copied) & 480~MB \\
Binaries copied, full run: the 62 referenced binaries not authenticated & 250~MB \\
Binaries copied, collection only: 2,853 not authenticated out of 44,446 read & 1.2~GB \\
Event logs (\opt{events}) & 170~MB \\
Hives, Prefetch, shortcuts, jump lists & 180 to 230~MB \\
\bottomrule
\end{tabular}
\end{center}

A machine used for a long time, or loaded with software, exceeds these
figures: the \path|$MFT| grows with the number of files, the unauthenticated
third-party binaries with the installed software. Allow \textbf{at least
5~GB free} for a collection only with \opt{binary}.

\wac{} checks the location \textbf{before its very first write} and refuses to
start:
\begin{itemize}
  \item if the working directory already holds files — an earlier collection
        would be analysed instead of this one;
  \item if the free space is clearly insufficient.
\end{itemize}
If space runs out while binaries are being collected, they are hashed without
being copied, and the investigation log says so.

\section{Before going on site}

\begin{itemize}
  \item Use a \textbf{new output folder} for each machine (option
        \opt{output}).
  \item Note the time of the intervention and the operator's identity: \wac{}
        records them too, but an independent record is useful.
  \item Keep in mind that plugging the stick in and running \wac{} leave
        unavoidable traces (chapter~\ref{chap:traces}): they are known and
        documented, so that they are not mistaken for the activity under
        study.
\end{itemize}

% ============================================================ 4
\chapter{Running a collection}

\section{Command line}

Open a command prompt \textbf{as administrator}, move to the collection medium,
then run:

\begin{lstlisting}
WAC.exe [--collect | --convert=folder] [--events] [--binary | --binary-all]
        [--threads=N] [--dump]
        [--output=folder] [--loglevel=N] [--debug]
WAC.exe --update-trust[=folder]
\end{lstlisting}

Every option is optional and disabled by default.

\section{Configuration file: \fichier{wac.yml}}

\textbf{Version.} \wac{} numbers its versions \texttt{<major>.<minor>.<bugfix>}:
major for a rework, minor for features added, bugfix for fixes. The version
is read in \fichier{WAC.exe}'s properties (Details tab), shown at start, and
recorded in \fichier{investigation.json} (\champ{Tool.Version}): a collection
says which version made it. This document describes version 1.4.0.

The last build is in \fichier{bin} of the repository: \fichier{WAC.exe},
the reference \fichier{wac.yml}, the scripts that verify a collection
(section~\ref{sec:verif}) and \fichier{SHA256SUMS} to check them all.
Copying \fichier{WAC.exe} and \fichier{wac.yml} to the collection key is
enough; the scripts go with the investigator.

Rather than typing the options at every collection, the procedure is written
once in \fichier{wac.yml}, next to \fichier{WAC.exe}: \wac{} applies it
first, and an option given on the command line overrides it.
\opt{write-config} writes the reference file, commented — the recommended
procedure: collection only, unverified binaries collected, event logs
included.

\begin{lstlisting}
convert: false          # true: collect and convert here; false: collect only
binary: unverified      # none | unverified | all
threads: 0              # analysis threads for the binaries, 0: automatic
output: output
log_level: 0
dump: false
language: fr            # fr | en: language of the console; the files stay in English
artefacts:              # one switch per artefact; all the hives = registry
  registry: true
  events: true
  scheduled_tasks: true
  services: true
  processes: true
  sessions: true
  network: true
  prefetch: true
  jump_lists: true
  recent_documents: true
  mft: true             # $MFT: every file, present or deleted, with its dates
  usn_journal: true     # $UsnJrnl:$J: creations, renames, deletions
  log_file: true        # $LogFile: NTFS transaction journal, raw copy only
  powershell_history: true  # PSReadLine: the commands typed in PowerShell
  setupapi: true        # setupapi.dev.log: device installations
  recycle_bin: true     # $I of the recycle bins: what, by whom, when
  hosts: true           # the hosts file: names diverted to an address
  error_reports: true   # error reports (WER): programs that crashed, with their path
  defender: true        # Defender: logs (MPLog), detection history, quarantine (copied encrypted)
  startup_folders: true # startup folders: what runs at logon
  zone_identifiers: true  # Zone.Identifier streams: where downloaded files came from
  security: true        # SECURITY hive: LSA secret names, cached logons, audit
  persistence: true     # registry persistence: IFEO, AppInit, Winlogon, LSA, COM…
  office_trust_records: true  # Office trusted documents, macros enabled
  typed_paths: true     # paths typed in Explorer
  word_wheel_query: true  # Explorer searches
  root_certificates: true # root certificates of the machine and users
  internet_settings: true # users' proxy
  firewall: true        # firewall rules and profiles
  wmi: true             # WMI repository read offline: subscriptions (filters, consumers, bindings)
  srum: true            # SRUDB.dat: network and resources per application and user, 30 to 60 days
  bits: true            # qmgr.db: BITS jobs, URLs, destinations, notification command lines
  fat_volumes: true     # FAT and exFAT volumes, the suspect's USB keys included: entries, deleted ones too
  installed_programs: true  # installed programs (Uninstall) of the machine and of each user
\end{lstlisting}

\begin{itemize}
  \item \textbf{Strict}: an unknown key, a value outside the ones allowed
    (\texttt{true}/\texttt{false} only, not \texttt{yes}), a duplicate key, a
    list: the collection is refused before it starts, with the line at fault.
    A key left out keeps its default.
  \item \textbf{Recorded}: the file's path, SHA-256 and content go into
    \fichier{investigation.json} (\champ{Tool.Configuration}) — the procedure
    applied is part of the evidence.
  \item An artefact switched off is not read at all, and its JSON file says so
    (\champ{CollectionStatus} \texttt{NotRequested}, with the key): never
    mistaken for a failure or an absence of traces.
  \item \opt{config=file} reads another file; \opt{full} collects and converts
    here although the file says \texttt{convert: false}.
  \item Without \fichier{wac.yml}, \wac{} behaves as before: full run, options
    as given.
\end{itemize}

\section{Options}

\begin{longtable}{@{}>{\ttfamily}L{3.2cm}L{11.2cm}@{}}
\toprule
\normalfont\bfseries Option & \bfseries Effect \\
\midrule
\endhead
-{}-collect & Collection only, on the examined machine: live observations and
  raw extraction into the sealed exhibit store; nothing is converted. See
  chapter~\ref{chap:separation}. \\
-{}-convert=folder & Conversion only, on an analysis workstation, of the
  collection made with \opt{collect} in that folder, after checking the seal
  and every fingerprint. See chapter~\ref{chap:separation}. \\
-{}-events & Extracts and parses the event logs (\fichier{.evtx}), and rebuilds
  each event's readable message from its provider's resources. Adds about
  250~MB of copies. \\
-{}-binary & Fingerprints (MD5, SHA-1, SHA-256) every file referenced by an
  artefact, by reading the disk directly. Referenced executables, libraries,
  drivers and scripts are also \textbf{collected} into the exhibit store. See
  section~\ref{sec:binary}. \\
-{}-binary-all & Like \opt{binary}, but \textbf{every} executable is collected,
  authenticated or not; the signature is checked all the same, and the manifest
  says for each whether it holds. For the cases where the content itself must
  be in the exhibit store. \\
-{}-threads=N & Number of threads analysing the executables of \opt{collect}
  \opt{binary} (PE, Authenticode, catalogs, fingerprints). By default, the
  processor's threads minus one, the disk reading keeping its own. The result
  does not depend on it: the manifest is the same. \\
-{}-update-trust\newline [=folder] & On the \textbf{analysis workstation}, \textbf{before} the
  collection, and alone: downloads from Microsoft the list of trusted roots
  (\fichier{authroot.stl}), that of disallowed certificates
  (\fichier{disallowedcert.stl}) and every root, checks their signatures and
  fingerprints, together with the vulnerable driver lists of Microsoft and
  LOLDrivers (not signed: authenticated by HTTPS alone) and the revocation
  lists (CRL) of the code signing authorities the CCADB lists, of Microsoft's
  own, and of those an earlier collection found missing
  (\fichier{wanted-crls.txt}), and writes them into \fichier{folder} (by default
  \fichier{trust}, next to \fichier{WAC.exe}), to be carried on the key with
  \fichier{WAC.exe}. The only mode that uses the network: the network library
  is loaded there alone. Refused with any other option. \\
-{}-dump & Adds the raw hexadecimal value of shellbag and \fichier{.lnk} shortcut
  items, for later decoding. \\
-{}-output=folder & Output folder, relative to the current directory. Default:
  \fichier{output}. \\
-{}-language=fr|en & Language of the console: steps, help, error messages.
  French by default; \texttt{language:} in \fichier{wac.yml} sets it too. The
  JSON files, \fichier{investigation.json} and the text log stay in English
  whatever the language: they are evidence read by tools, and a collection
  gives the same files whoever the operator. \\
-{}-loglevel=N & Enables a detailed text log in the current directory, named
  after the program run (\fichier{WAC.exe.log} for \fichier{WAC.exe}): 0~none,
  1~per artefact type, 2~per artefact, 3~per sub-function (debugging). \\
-{}-debug & Prints the NTFS reader's trace on standard error (path resolution,
  index blocks, cluster runs). For diagnosis only. \\
-{}-help, /? & Shows the help. \\
\bottomrule
\end{longtable}

\section{Examples}

Standard collection:
\begin{lstlisting}
E:\> WAC.exe --output=accounting-pc-01
\end{lstlisting}

Full collection, with logs, binaries and a detailed log:
\begin{lstlisting}
E:\> WAC.exe --events --binary --output=accounting-pc-01 --loglevel=2
\end{lstlisting}

\section{The \opt{binary} option}\label{sec:binary}

A fingerprint is enough to query a public database (VirusTotal, for instance)
\textbf{without sending it the file}. But it says nothing about a binary nobody
knows — precisely the one that matters to the investigation — and a binary left
behind may be gone by the time a detection comes in. This is why \opt{binary}
also copies potential payloads:

\begin{itemize}
  \item \textbf{collected}: \fichier{exe}, \fichier{dll}, \fichier{sys},
        \fichier{ocx}, \fichier{cpl}, \fichier{scr}, \fichier{drv},
        \fichier{efi}, \fichier{com}, \fichier{msi} and the scripts
        \fichier{ps1}, \fichier{psm1}, \fichier{bat}, \fichier{cmd},
        \fichier{vbs}, \fichier{vbe}, \fichier{js}, \fichier{jse},
        \fichier{wsf}, \fichier{wsh}, \fichier{hta};
  \item \textbf{collected as well}: Office documents \textbf{able to carry
        macros}, an intrusion vector as common as an executable — legacy binary
        formats (\fichier{doc}, \fichier{xls}, \fichier{ppt}, \fichier{pps},
        \fichier{dot}, \fichier{xlt}, \fichier{pot}), ``m'' formats
        (\fichier{docm}, \fichier{dotm}, \fichier{xlsm}, \fichier{xltm},
        \fichier{pptm}, \fichier{potm}, \fichier{ppsm}), \fichier{xlsb}, add-ins
        (\fichier{xla}, \fichier{xlam}, \fichier{xll}, \fichier{wll},
        \fichier{ppa}, \fichier{ppam}), Publisher (\fichier{pub}), Visio
        (\fichier{vsd}, \fichier{vsdm}, \fichier{vstm}, \fichier{vssm}) and
        Access (\fichier{mdb}, \fichier{accdb}, \fichier{accde}). They appear in
        the traces as targets of shortcuts or jump lists, or among the files
        loaded by Word or Excel in their Prefetch;
  \item \textbf{only hashed}: other documents and data files, including
        \fichier{docx}, \fichier{xlsx} and \fichier{pptx}, which cannot contain
        macros. Copying them would turn the collection into a copy of the
        user's files.
\end{itemize}

The referenced files come from processes, services (binary and DLL), scheduled
tasks, files loaded by programs (Prefetch), Shimcache, Amcache, and targets of
shortcuts and jump lists. Each file is read only once, however many artefacts
cite it, and \textbf{identical content is stored only once}: the other paths
appear in the manifest as exhibits sharing that content
(\champ{SharedExhibit}).

\subsection*{Authentic Microsoft binaries: hashed, not collected}

Most referenced binaries are Windows components, identical on every machine of
the same build. Their origin can be proven on the machine itself, with no list
to carry: Windows' \textbf{catalogs}, signed by Microsoft, list the fingerprint
of every system file. \wac{} reads those of \path|System32\CatRoot|, those of
the components (\path|WinSxS\Catalogs|) and those of the servicing packages
(\path|servicing\Packages|), over 13,000 on Windows~11. Individually signed
binaries (Edge, OneDrive, Office, .NET) carry an \textbf{embedded signature}. A
binary authenticated this way is hashed without being collected; every other one
is collected — including a system binary replaced by an attacker, whose
fingerprint no longer matches.

This check leaves \textbf{no trace}: catalogs and binaries are read directly
from the disk, and the cryptographic verification is done in memory by \wac{},
up to Microsoft root certificates embedded in the executable. No service is
solicited, the machine's certificate store is not consulted, and revocation is
not checked online.

Only signers Microsoft uses for its own code are accepted (``Microsoft
Windows'', ``Microsoft Corporation'', ``Microsoft Windows Publisher'' and
``.NET''). \textbf{Third-party}
drivers certified by Microsoft (signers ``Hardware Compatibility Publisher'' and
``Early Launch Anti-malware Publisher'') are still collected: a vulnerable but
signed driver is a classic attack path.

\subsection*{Store applications: verified through their package}

The files of a Store application are not signed one by one: their
\textbf{package} is. Its folder holds the signature (\fichier{AppxSignature.p7x})
and the block map (\fichier{AppxBlockMap.xml}), which gives the SHA-256 of every
64~KiB block of every file. \wac{} verifies the signature, which must reach a
Microsoft root (the Store signs the packages it distributes), checks that it
does cover that map, then compares every block read with the map. A matching
file is hashed without being collected; a modified file, or one added to the
folder after installation, is collected, and so is every file of a package
signed outside the Microsoft chain.

The catalogs, package signatures and block maps that justified not collecting a
file go into the exhibit store, so that a third party can redo the check. In
the artefact files, the \champ{Signature} field (or \champ{ServiceDllSignature},
\champ{SignatureTarget}\ldots) states the authentication used, for instance
``\texttt{Microsoft (catalog C:\textbackslash{}Windows\textbackslash{}\ldots.cat)}'',
``\texttt{Microsoft (embedded signature)}'' or
``\texttt{Package (package C:\textbackslash{}Program Files\textbackslash{}WindowsApps\textbackslash{}\ldots)}''.

\textbf{Scripts} are authenticated the same way: by Windows' catalogs and, for
PowerShell (\fichier{ps1}, \fichier{psm1}, \fichier{psd1},
\fichier{ps1xml}\ldots), by their embedded signature block
(\texttt{\# SIG \# Begin signature block}). A Windows Script Host script
(\fichier{vbs}, \fichier{js}, \fichier{wsf}) signed only within its text, and
not by a catalog, is still collected.

On the test machine, this check brings collection down from 2,131 binaries
(3.3~GB) to 92 (335~MB), without lengthening the collection; the scripts still
collected are the unsigned ones.

Since reading is direct, collecting leaves no more trace than hashing — it only
takes space on the medium.

\section{What the screen shows}

\wac{} displays each step, followed by \texttt{OK} or an error message. A failed
step does not stop the collection: whatever is accessible is gathered, and
whatever is missing is recorded.

\begin{lstlisting}
[PREREQUISITE VERIFICATION]
 - Check OS >= Windows 10 : OK
 - Check administrator rights : OK
 - Checking the collection medium : OK
[SEARCHING FOR ARTIFACTS IN WINDOWS API]
 - Extraction of SESSIONS : OK
 - Extraction of PROCESS : OK
 - Extraction of NETWORK STATE : OK
[RAW EXTRACTION]
 - Extracting system hives (raw NTFS) : OK
 - loading the HKLM\SYSTEM key : OK
 - loading the HKLM\SOFTWARE key : OK
 - Listing USER PROFILES (SOFTWARE hive) : OK
 - Extracting user hives (raw NTFS) : OK
 - Extracting file artefacts (raw NTFS) : OK
[SEARCHING FOR ARTIFACTS IN THE REGISTRY]
 - Extracting USERASSIST Registry Keys : OK
   ...
[SEARCHING FOR ARTIFACTS IN FILES]
 - Extracting PREFETCHS : OK
   ...
[SEARCHING FOR ARTIFACTS IN EVENT LOGS]
 - Extraction of EVENTS : OK
[EXHIBIT STORE]
 - Copying collected binaries to the working directory : OK
 - Sealing the exhibit store (manifest + SHA-256) : OK
[INVESTIGATION LOG]
 - Writing investigation.json : OK
END, Time elapsed : 309 s
\end{lstlisting}

This is the display with \opt{language=en}, or \texttt{language: en} in
\fichier{wac.yml}: by default the console is in French.

Sealing the exhibit store is the \textbf{last} operation on it: every exhibit,
including those added while the event logs are processed, is covered by the
manifest.

\section{Duration}

The duration depends mostly on the hardware (USB~2 or~3, disk, processor) and
the options. On the test virtual machine (NVMe disk): about 70~s with
\opt{events}, about 310~s with \opt{events} and \opt{binary}. On a USB stick,
writing dominates.

% ============================================================ 4 bis
\chapter{Collect here, convert elsewhere}\label{chap:separation}

\section{Why separate}

A run of \wac{} does two different things:
\begin{enumerate}
  \item \textbf{the collection}: observing what only a running system shows
        (processes, sessions, service states, clock), reading the disk
        directly, copying the evidence into a sealed exhibit store. It
        \textbf{needs} the examined machine;
  \item \textbf{the conversion}: parsing thousands of copied files (hives,
        event logs, Prefetch, shortcuts\ldots) into JSON. It only needs the
        exhibit store, \textbf{not} the machine.
\end{enumerate}

Without a mode option, both run on the examined machine. The conversion adds
work there that nothing requires: minutes of processing in \wac{}'s memory,
which Windows may page out to \fichier{pagefile.sys}, on the very disk under
examination; a process present for longer, seen and logged by the machine's
security software; if it crashes, a Windows error report written on the
machine. It reads neither the machine's files nor its registry, only the copies
on the collection medium, but it runs there.

That is exactly what good practice asks to avoid:
\begin{itemize}
  \item \textbf{ISO/IEC~27037} covers the identification, collection and
        preservation of digital evidence, \textbf{ISO/IEC~27042} its analysis:
        two distinct stages;
  \item \textbf{ACPO principle~1}: no action should change the data of the
        examined machine; every operation run there must be necessary;
  \item \textbf{ACPO principle~3}: an independent third party must be able to
        repeat the process and reach the same result.
\end{itemize}

Separating the two stages brings three things:
\begin{itemize}
  \item \textbf{fewer traces} on the examined machine: only the collection runs
        there;
  \item \textbf{evidence checked before any analysis}: \opt{convert} first
        verifies the seal and the SHA-256 of every exhibit, and refuses a
        retouched collection;
  \item \textbf{a reproducible analysis}: two conversions of one collection
        give byte-identical JSON files, and the analysis can be redone later,
        with a corrected version of \wac{} or by a third party, without going
        back to the machine.
\end{itemize}

\section{Which mode}

\begin{itemize}
  \item \textbf{By default: \opt{collect} on the machine, \opt{convert} on the
        analysis workstation.} It is the procedure that leaves the fewest traces
        on the examined machine, keeps the evidence sealed, and allows the
        analysis to be redone and checked.
  \item \textbf{Full run} (no mode option): for specific cases, to be stated in
        the report — a triage whose results are needed on the spot with no
        analysis workstation, a machine to be released within minutes. Its
        conversion reads only the copies on the collection medium, never the
        machine's files or registry, but it runs on it: minutes of processing
        in memory (which Windows may page to \fichier{pagefile.sys}), a longer
        process, a Windows error report if it crashes.
  \item \textbf{Scope of \opt{binary}}: in a full run, only the files
        \textbf{cited by the artefacts} are read; a binary on the disk that no
        artefact cites is not examined. With \opt{collect}, it is \textbf{every
        executable of every fixed NTFS volume}, cited or not. The collection
        alone takes longer for that reason (8~min~30 against 3~min~30 on the
        test machine): that time is reading, which writes nothing to the
        machine.
  \item \textbf{\opt{binary-all}}: when the content of every executable must be
        in the exhibit store (malware analysis, request for the files
        themselves), for about 20~GB.
\end{itemize}

\section{The two commands}

On the examined machine, as administrator:
\begin{lstlisting}
E:\> WAC.exe --collect --events --binary --output=accounting-pc-01
\end{lstlisting}

On the analysis workstation, with no particular rights, with the same
\opt{events} and \opt{binary} options as the collection:
\begin{lstlisting}
D:\> WAC.exe --convert=E:\accounting-pc-01 --events --binary
\end{lstlisting}

The JSON files are written into the collection's folder, next to the exhibit
store. \opt{convert} excludes \opt{collect} and \opt{output}.

\section{What \opt{collect} takes}

Since the conversion can no longer ask the machine anything, the collection
takes \textbf{everything it will read}:
\begin{itemize}
  \item the \textbf{live observations}, recorded as read in
        \fichier{exhibits\textbackslash live} and sealed like any exhibit:
        processes, sessions, service states, clock and boot, and the machine's
        settings (system drive, code page, time zone, volumes and their serial
        numbers);
  \item the hives and the file artefacts, as usual;
  \item with \opt{events}, the resource files of \textbf{every} declared event
        provider (binary, parameter file, \fichier{.mui} satellites): which of
        them are needed is only known at conversion. About 650~MB on
        Windows~11;
  \item with \opt{binary}, \textbf{every} executable, library, driver, script
        and macro document of \textbf{every fixed NTFS volume} — a profile on
        \fichier{D:} included —, read directly and verified in memory, as in a
        full run (section~\ref{sec:binary}):
        \begin{itemize}
          \item authenticated (catalog, embedded signature, Store package), or a
                third-party binary whose signature holds against the trust set of
                the key (\opt{update-trust}: chain to a root Microsoft trusts for
                code signing, valid at the signing time, not revoked, not a
                vulnerable driver): \textbf{not copied}; the manifest keeps its
                path, its Authenticode digests and the verdict
                (\champ{ContentStored} \texttt{false}); without a usable trust
                set, every third-party binary is copied, as the investigation log
                says;
          \item not authenticated: \textbf{copied}, with its three fingerprints
                MD5, SHA-1, SHA-256, its Authenticode digest and the reason;
          \item Windows differential file (\texttt{PA30}/\texttt{PA31}, the
                \fichier{r} and \fichier{f} folders of \fichier{WinSxS}): not an
                executable; hashed, not copied.
        \end{itemize}
        A content is stored once, and the other names of a hard link
        (\fichier{System32} and \fichier{WinSxS}) are recorded in the manifest
        like the first. Removable volumes, the collection medium among them,
        and the memory files (\fichier{pagefile.sys}, \fichier{swapfile.sys},
        \fichier{hiberfil.sys}) are left out.
\end{itemize}

\begin{attention}
With \opt{binary}, the collection reads every volume: on the test virtual
machine (Windows~11 without added software), 44,446~executables read, of which
41,593 authenticated (or third-party with a valid signature) not copied and
2,853 copied (1.2~GB); the whole collection takes 8~minutes~30 and the
exhibit store, event logs and NTFS metadata included, weighs 2.7~GB (see
section~\ref{sec:place}). If space runs short, the remaining files are hashed without being
copied, and the manifest says so.
\end{attention}

No artefact JSON file is written: the folder holds only the sealed exhibit
store, the working directory and \fichier{investigation.json}.

\section{What \opt{convert} does}

\begin{enumerate}
  \item \textbf{Check}, before any read: the seal must carry the manifest's
        real fingerprint, and every exhibit its SHA-256 fingerprint. A single
        altered exhibit, and the conversion is \textbf{refused};
        \fichier{conversion.json} gives the reason.
  \item \textbf{Rebuild of the working directory} from the exhibit store: the
        former one is deleted rather than trusted, the copies are verified
        against the manifest and the hives repaired, as at collection time. The
        executables and the providers' resources are read in the exhibit store
        itself, without a copy.
  \item \textbf{Conversion}, from the exhibits alone: nothing is read from the
        converting machine — neither its disk, nor its registry, nor its time
        zone.
\end{enumerate}

The exhibit store is never modified and its seal remains the collection's.
\fichier{investigation.json} remains the collection's log;
\fichier{conversion.json} is the conversion's (operator, workstation,
operations). Two conversions of one collection give identical JSON files,
which the test harness checks on every cycle.

\section{Limitations}

\begin{itemize}
  \item \textbf{Macro documents} (\fichier{doc}, \fichier{docm},
        \fichier{xls}, \fichier{xlsm}, \fichier{ppt}\ldots, list in
        section~\ref{sec:binary}) are handled like executables: read on every
        fixed volume and \textbf{collected}, unless they are, byte for byte, a
        file shipped by Microsoft and listed in a catalog — which a document
        written or modified by anyone cannot be.
  \item A cited file that can carry neither code nor macro (\fichier{docx},
        \fichier{xlsx}, \fichier{pdf}, data) is not collected: at conversion it
        therefore has \textbf{no fingerprint}, whereas the full run computes it
        on the disk.
  \item An authenticated executable is not copied: at conversion it has its
        Authenticode digest and its verdict, but no MD5, SHA-1 or SHA-256 —
        those are computed only for what is collected.
  \item An executable on a removable volume is not collected by
        \opt{collect} \opt{binary}.
  \item The \opt{events} and \opt{binary} options of the conversion must be
        those of the collection: what the collection did not take cannot be
        converted.
\end{itemize}

% ============================================================ 5
\chapter{The results}

\section{Layout}

\begin{lstlisting}
<output>/
  exhibits/            exhibit store: raw copies, never rewritten
    MANIFEST.json     identifies each exhibit and the collection
    MANIFEST.sha256   the manifest's fingerprint: its seal
    Windows/...        the exhibits, under their original paths
    Users/...
  working/             working copies (replayed hives)
  investigation.json   the collection log
  *.json               the artefacts
\end{lstlisting}

A file from a volume other than the system one is stored under
\path|_volume_D\...|; the JSON always gives the original path, with its real
drive letter.

\section{The exhibit store manifest}

\fichier{exhibits/MANIFEST.json} identifies the collection and each exhibit.

\subsection{Header}

\begin{description}[style=nextline,leftmargin=1.2em]
  \item[\champ{Tool}] tool name, command line, build date.
  \item[\champ{Host}] machine name, system drive, time zone of the suspect (read
        from its \fichier{SYSTEM} hive) and of the collecting machine.
  \item[\champ{Operator}] account that ran the collection, and whether it was
        elevated.
  \item[\champ{Custody}] procedure statement, start of extraction, time the
        manifest was written, number of exhibits and failures, total size, hash
        algorithms, volumes read (letter, serial number, file system).
\end{description}

\subsection{Each exhibit}

\begin{longtable}{@{}>{\ttfamily}L{4.2cm}L{10.2cm}@{}}
\toprule
\normalfont\bfseries Field & \bfseries Meaning \\
\midrule
\endhead
SourcePath & Original path, with drive letter. \\
ExhibitPath & Location of the raw copy, relative to the output folder. \\
Method & Collection method (direct NTFS reading, volume, purpose). \\
Result & \texttt{OK}, or the error code and message. Failures are listed in the
  manifest: a missing exhibit would read as one never looked for. \\
MD5, SHA1, SHA256 & The three fingerprints, computed while writing: they bear
  on what was read from the volume. \\
Bytes & Extracted size. \\
DeclaredBytes & Size declared by NTFS, present only when it differs from the
  extracted size (truncated extraction). \\
ValidDataBytes & Valid data length of a preallocated file (event logs): beyond
  it, the exhibit holds zeros by definition. \\
SharedExhibit & \texttt{true} when the content, identical byte for byte, is
  already stored under another exhibit: \champ{ExhibitPath} points to it. The
  source remains an exhibit in its own right, with its own path, \path|$MFT|
  entry and timestamps. \\
MftEntry & NTFS record number: identifies the file independently of its name. \\
Extracted(Utc) & When this exhibit was extracted. \\
SourceCreated, SourceModified, SourceMftModified, SourceAccessed & The four NTFS
  timestamps of the original file (\path|$STANDARD_INFORMATION|), in UTC and in
  the suspect's local time. They, not the dates of the copies, date the file. \\
SourceFileName\newline Created, Modified, MftModified, Accessed &
  The four timestamps of the \path|$FILE_NAME| attribute of the name read. NTFS
  keeps them apart, and a change of the dates by a program leaves them
  untouched: a \champ{SourceCreated} earlier than \champ{SourceFileNameCreated}
  is the classic sign of dates forged after the fact (\emph{timestomping}). \\
ContentStored & \texttt{false}: authenticated binary read but not copied
  (\opt{collect} \opt{binary}); it has no \champ{ExhibitPath}. \\
AuthenticodeSHA1, AuthenticodeSHA256 & Authenticode digests of a PE, those its
  signature and the catalogs cover; the only digests of an authenticated binary
  that was not copied. \\
SignatureVerified & Outcome of an executable's signature check: \texttt{true}
  with \champ{Signature} (what authenticates it), \texttt{false} with
  \champ{SignatureReason} (why not). \\
PeTimeDateStamp, PeSizeOfImage, SymbolServerKey & Build of a PE and the key
  (\texttt{name/TIMESTAMPsize/name}) under which Microsoft's symbol server keeps
  it: an authentic Microsoft binary that was not copied can be fetched again
  from it, identical, and its fingerprint checked. \\
\bottomrule
\end{longtable}

\begin{note}[Why three fingerprints]
MD5 collisions have been producible at will since 2008, and SHA-1 ones since
2017: a single fingerprint no longer settles an identity dispute. Three
independent ones do.
\end{note}

\section{Verifying a collection}\label{sec:verif}

\subsection{The seal}

The manifest cannot hold its own fingerprint: it is in
\fichier{MANIFEST.sha256}. On Linux:

\begin{lstlisting}
$ cd output/exhibits
$ sha256sum -c MANIFEST.sha256
MANIFEST.json: OK
\end{lstlisting}

\subsection{Each exhibit}

Two scripts, delivered with \fichier{WAC.exe} in \fichier{bin}, check the
seal, then the SHA-256 of every exhibit: \fichier{verify-exhibits.ps1} on
Windows, \fichier{verify-exhibits.sh} on Linux. Their French counterparts,
\fichier{verifier-consigne.ps1} and \fichier{verifier-consigne.sh}, make the
same checks with French messages.

\textbf{Verify from a Linux machine.} This is the recommended way: the
collection may hold what infected the examined machine, and a Linux
investigation machine is a technological break — a Windows executable or macro
document of the collection cannot run there. The script needs only
\fichier{bash}, the coreutils and \fichier{python3} (which reads the JSON);
the fingerprints are computed by \fichier{sha256sum}, not by \wac{}:

\begin{lstlisting}
$ bash verify-exhibits.sh /media/usb/output
Seal verified (CEB489A143AE29EB9E3BC836F3D565B744F036AE1F3A72A3447408F31F8A5B67)
1617 exhibit(s) verified, 0 mismatch(es)
0 binarie(s) authenticated and fingerprinted only (no copy to verify)
\end{lstlisting}

The manifest comes from the examined machine: both scripts refuse an exhibit
path that leaves the collection (\texttt{..}, absolute path, and for the
\fichier{.sh} a symbolic link), and count it as a mismatch without reading it.

\textbf{From Windows:}

\begin{lstlisting}
PS> powershell -ExecutionPolicy Bypass -File verify-exhibits.ps1 -Output E:\output
Seal verified (CEB489A143AE29EB9E3BC836F3D565B744F036AE1F3A72A3447408F31F8A5B67)
1617 exhibit(s) verified, 0 mismatch(es)
\end{lstlisting}

Their exit code is~0 if everything matches, 1~otherwise.

\lstinputlisting[language={},caption={}]{verify-exhibits.sh}

\lstinputlisting[language={},caption={}]{verify-exhibits.ps1}

\section{The investigation log}

\fichier{investigation.json} traces the collection: tool and command line,
machine and time zones, operator, start, end and duration, then the
\textbf{ordered} list of operations. For each one:

\begin{description}[style=nextline,leftmargin=1.2em]
  \item[\champ{Seq}, \champ{TimestampUtc}, \champ{TimestampLocal}] rank and time;
  \item[\champ{Operation}, \champ{Target}] what was done, and on what;
  \item[\champ{Result}] the outcome;
  \item[\champ{Footprint}] the trace this operation is expected to leave on the
        examined machine.
\end{description}

This log makes it possible to compare, collection by collection, the expected
traces with those actually observed. It also flags a difference between the
suspect's time zone and the collecting machine's.

\section{The artefact files}

One file per artefact. The fields described are those emitted; a field without
a value is not emitted (see chapter~\ref{chap:lecture}).

\subsection{System and accounts}

\begin{description}[style=nextline,leftmargin=1.2em]
  \item[\fichier{OperatingSystem.json}] machine name and domain, Windows edition
        and version, installation date, current time, \textbf{boot time},
        uptime, time zone. \champ{BootTimeSource} says where the boot time comes
        from (kernel, or fallback estimate); \champ{ClockAdjustedSinceBootMs}
        gives the total clock corrections since boot — a few seconds is routine
        resynchronisation, a large value means a clock changed.
  \item[\fichier{security.json}] \fichier{SECURITY} hive: \textbf{names} of the
        LSA secrets and the dates of their values (\texttt{DefaultPassword}
        reveals an automatic logon's stored password, \texttt{\_SC\_<service>}
        a service running under an account), number of cached domain logon
        slots and how many are used, audit policy (\champ{PolAdtEv}, as bytes:
        its layout is not published). \textbf{No secret's value} is read, nor
        any cached credential. The \path|config\RegBack| copies of the hives
        are collected as they are.
  \item[\fichier{users.json}] local accounts, read from the \fichier{SAM} hive:
        SID and RID, state, flags, logon and failure counts, dates of last
        logon, last failure, password change and account change, profile.
  \item[\fichier{Sessions.json}] sessions open at collection time: user, logon
        type, authentication package, start time.
  \item[\fichier{processes.json}] running processes: name, identifiers (and
        parent), owner, session, loaded modules and, with \opt{binary}, the
        executable's fingerprints.
  \item[\fichier{network.json}] network state at collection time:
        \champ{Connections} (TCP and UDP, IPv4 and IPv6: addresses, ports,
        state, owning process and its name, creation time local and UTC),
        \champ{Neighbours} (ARP and NDP caches: the machines of the local network
        spoken with lately), \champ{Interfaces} (addresses, gateways, DNS
        servers, DHCP) and \champ{Routes}. \champ{TablesNotRead} names a table
        Windows refused to give, which is not an empty table.
  \item[\fichier{services.json}] services and drivers, configuration read from
        the hive: type, start mode, account, binary and hosted DLL, failure
        command, dependencies, last write of the key; current state; with
        \opt{binary}, fingerprints of the binary and of the DLL.
  \item[\fichier{ScheduledTasks.json}] scheduled tasks, read from their XML
        definitions: author, run-as account, state, last run and its result,
        triggers and actions (with \opt{binary}, fingerprints of the command).
\end{description}

\subsection{Program execution}

\begin{description}[style=nextline,leftmargin=1.2em]
  \item[\fichier{prefetchs.json}] Prefetch files: program, run count and times,
        volumes, and list of loaded files with their \path|$MFT| reference (and
        their fingerprints with \opt{binary}).
  \item[\fichier{amcache\_application\_files.json},
        \fichier{amcache\_applications.json},
        \fichier{amcache\_inventory.json}] Amcache inventory, table by table
        and value by value, with the last write time of each entry: the
        executables (path, publisher, version, size, link date, \champ{Sha1}),
        the installed applications, and in the third file the drivers (with
        their \champ{Sha1}), devices, containers, driver packages and
        shortcuts, each with its \champ{Table}. \textbf{Beware}: for a file
        larger than 30~MiB, Amcache hashes only the first 31,457,280 bytes;
        \champ{Sha1CoversBytes} says so, and that SHA-1 does not compare with the
        whole file's.
  \item[\fichier{shimcache.json}] compatibility cache: path and modification date
        of the file, position in the cache (\champ{CachePosition}, 0 for the most
        recent). An entry proves the binary was present, not that it ran. The
        execution flag other tools publish is \emph{not} collected: measured
        against Prefetch, its ``yes'' is nearly always right, but its ``no'' is
        wrong for 9 executed binaries out of 10 — it would lead to conclude that
        a program never ran.
  \item[\fichier{bams.json}] last execution per program and per user
        (Background Activity Moderator). The hives being read from the disk,
        an execution of the last minute or so may be missing: Windows keeps a
        registry change in memory before writing it (56 to 82~s measured on
        the test machine, which \texttt{RegFlushKey} did not bring forward),
        and \wac{} does not force that write, which would modify the examined
        disk.
  \item[\fichier{userassists.json}] programs launched from the user interface:
        count, number and duration in the foreground (\champ{FocusTimeMs}), last
        execution.
  \item[\fichier{muicache.json}] display names of executed programs.
  \item[\fichier{run.json}] programs started at boot or logon (Run and RunOnce
        keys), with the key's last write: in SOFTWARE (the whole machine), in
        the DEFAULT hive (the system account) and in each profile's NTUSER.DAT
        (\champ{Hive}). \champ{Type} \texttt{REG\_EXPAND\_SZ}: the command
        holds variables (\texttt{\%SystemRoot\%}\ldots), given as written.
  \item[\fichier{persistence.json}] registry persistence beyond the Run keys:
        IFEO (\texttt{Debugger}, \texttt{GlobalFlag}), SilentProcessExit,
        AppInit\_DLLs and AppCertDlls, Winlogon (\texttt{Shell},
        \texttt{Userinit}, \texttt{Taskman}, \texttt{Notify}), LSA packages,
        print monitors, \texttt{BootExecute}, Active Setup, NetSh helpers, and
        COM hijacking in each user's classes (UsrClass.dat). For each value:
        mechanism (\champ{Mechanism}), hive, key, name, type, value as written,
        the key's last write, and with \opt{binary} the fingerprints of the
        file it names. Every value is given, Windows' own included: telling the
        added from the default is the analysis's work. BITS jobs are not read
        yet.
  \item[\fichier{office\_trust\_records.json}] Office documents a user trusted
        (``Enable content''): application, version, document, date;
        \champ{MacrosEnabled} when macros were enabled — the door of most macro
        infections.
  \item[\fichier{typed\_paths.json}, \fichier{word\_wheel\_query.json}] paths
        typed in Explorer's address bar and terms searched there, per user and
        position (\champ{MruPosition}, 0 = the most recent, as in
        \fichier{shellbags.json}).
  \item[\fichier{root\_certificates.json}] certificates of the machine's and each
        user's root stores: subject, issuer, validity; a root added makes Windows
        trust whatever it signs (TLS interception, signed malware).
        \champ{ThumbprintMismatch} flags a certificate whose SHA-1 is not its
        key's name.
  \item[\fichier{internet\_settings.json}] each user's proxy and automatic
        configuration script.
  \item[\fichier{firewall.json}] firewall rules (name, action, direction, ports,
        program; every field in \champ{Fields}) and the domain, private and
        public profile settings.
  \item[\fichier{installed\_programs.json}] installed programs, read in the
        \texttt{Uninstall} keys of the machine (64- and 32-bit) and of each user
        (\champ{Scope}): all their values (name, publisher, version, folder,
        uninstall command, \champ{InstallDate} --- the day only, as written) and
        the key's write time. A remote access tool installed, then uninstalled,
        often leaves its trace there.
  \item[\fichier{fat\_volumes.json}] FAT12, FAT16, FAT32 and exFAT volumes, the
        suspect's USB keys and cards included (the collection's volume
        excepted): every entry, present or \textbf{deleted} (\champ{Deleted}),
        its path, 8.3 name, size, attributes and dates. A name deleted on FAT
        keeps its long name unless overwritten; otherwise its 8.3 name, first
        letter lost (\champ{NameIncomplete}). Each date is given in UTC and in
        local time derived from UTC, as in every artefact: one instant reads
        the same everywhere. FAT stores the local time Windows computed when it
        wrote it: a modification earlier than the creation is a date kept by a
        copy, written at the copy, and converted with the offset of that moment
        (\champ{ModifiedWrittenAt} \texttt{Copy}, \champ{ModifiedOffsetApplied},
        \champ{ModifiedAsStored}) — otherwise a winter file copied in summer
        would be an hour off. exFAT keeps its offset (\champ{ModifiedOffset}\ldots),
        the UTC date is exact. The collection copies only the volume's metadata
        (boot sector, FAT, directories), not the files' content; with
        \opt{binary}, a program of these volumes cited by an artefact (run
        from a key, for instance) is fingerprinted and collected as on NTFS.
  \item[\fichier{srum.json}] resource usage (SRUM, \path|sru\SRUDB.dat|), one
        entry per row of every table: bytes sent and received per application
        and per interface (an exfiltration shows there, with the program), CPU
        cycles, bytes read and written, connections; 30 to 60 days.
        \champ{Table} (the provider's GUID), \champ{Provider} (documented name:
        \texttt{NetworkUsage}, \texttt{ApplicationResourceUsage}\ldots), then the
        columns by name (letters and digits only), and \champ{App},
        \champ{UserSID}, \champ{UserName} resolved through
        \texttt{SruDbIdMapTable}. The database is copied live: the last hour is
        in its logs \path|SRU*.log|, copied into the exhibit store for a replay
        (\texttt{esentutl /r}) on a copy.
  \item[\fichier{bits\_jobs.json}] BITS jobs (\path|qmgr.db|): name,
        description, type, state, owner, \textbf{notification command line}
        (run when the transfer ends or fails: a persistence), delays, dates,
        and each file: source URL, destination, temporary file, sizes. The
        whole blob is given in hexadecimal: BITS~10 keeps more there (the
        creating program) at places not described.
  \item[\fichier{wmi\_persistence.json}] WMI event subscriptions, read
        \textbf{offline} in the CIM repository copied raw
        (\path|wbem\Repository|): no WMI query, nothing run on the machine. In
        \emph{every} namespace, each filter (WQL query), consumer (command
        line, VBScript or JScript script, log, mail, and any class an intruder
        defined), binding and timer, with all its properties
        (\champ{Properties}: name and value; \champ{CreatorSID} in text). A
        persistence invisible in the registry. The two dates at the start of
        an instance are not emitted: they are not creation dates.
  \item[\fichier{startup\_folders.json}] content of the startup folders, of all
        users (\champ{Scope} \texttt{AllUsers}) and of each profile
        (\texttt{User}, with \champ{SID}): every file, its dates and
        fingerprints; a shortcut is decoded (\champ{Shortcut}, as in
        \fichier{recentdocs.json}).
  \item[\fichier{powershell\_history.json}] PSReadLine history of each profile
        and each host (\champ{Host}: \texttt{ConsoleHost}, the ISE, VS~Code):
        the commands in the order they were typed, a command on several lines
        joined. No date per command: \champ{SourceModified} dates the last one.
  \item[\fichier{error\_reports.json}] Windows error reports
        (\fichier{Report.wer}) of the machine and of each profile: program, path,
        instant (\champ{EventTimeUtc}), problem signature (faulting module,
        version), every value of the report in \champ{Values}.
  \item[\fichier{defender\_log.json}] Microsoft Defender's support logs
        (\fichier{MPLog-*.log}): every line as written, dated when it starts with
        a UTC date. Run-time scans, detections, exclusions, updates.
  \item[\fichier{defender\_detections.json}] one entry per DetectionHistory
        file, which survives the clearing of the event log: threat, resource
        (file, key), account, process, initial detection date, quarantine
        entry, and the named \champ{ThreatTracking} properties (fingerprints,
        size, start time). The format is undocumented: only the fields whose
        position was confirmed against \texttt{MSFT\_MpThreatDetection} are
        named; the others are given as read (\champ{UnnamedFields}: offset,
        type, value), without a meaning made up.
  \item[\fichier{defender\_quarantine.json}] quarantine entries: threat,
        instant of the quarantine, original path of the removed file, its
        creation and modification dates, attributes, size, and the
        fingerprint of its content (\champ{ContentSha256}). The quarantine is
        copied \textbf{encrypted}, as Defender keeps it: nothing malicious is
        written in clear, on the collection medium or on the analysis
        workstation. The content is decrypted in memory only, for its
        fingerprint.
  \item[\fichier{hosts.json}] entries of the \fichier{hosts} file: address and
        names, line number; an untouched Windows has none.
\end{description}

\subsection{Files and browsing}

\begin{description}[style=nextline,leftmargin=1.2em]
  \item[\fichier{mft.json}] every file and directory of every fixed NTFS
        volume (\champ{Volume}: \texttt{C:}, \texttt{G:}\ldots; a record
        number only means something within its volume),
        present or \textbf{deleted} (\champ{Deleted}) as long as NTFS has not
        reused its record: rebuilt path, size, attributes, and two sets of four
        UTC dates — \champ{CreatedUtc}\ldots{} (the
        \path|$STANDARD_INFORMATION| attribute, what Explorer shows and what a
        timestomping tool rewrites) and \champ{FileNameCreatedUtc}\ldots{} (the
        \path|$FILE_NAME| attribute, written by NTFS alone). A
        \path|$STANDARD_INFORMATION| date earlier than the \path|$FILE_NAME|
        one betrays timestomping. \champ{PathIncomplete} marks a path one of
        whose parent directories was reused since: the path is not made up.
        A file with hard links (\champ{HardLinks}) is one record under several
        names — \path|System32\crypt32.dll| and its original under
        \path|WinSxS|: \champ{HardLinkPaths} gives its other paths, each one
        searchable.
  \item[\fichier{usn\_journal.json}] the change journal of every fixed NTFS
        volume (\path|$UsnJrnl:$J|, \champ{Volume}; the \champ{Usn} numbers
        rise within a volume): every creation, write, rename and
        \textbf{deletion}, dated, with its reasons (\champ{Reasons}, an unknown
        bit kept in hexadecimal) and the path rebuilt through the \path|$MFT|.
        It keeps the trace of files dropped then erased, which the \path|$MFT|
        has already forgotten.
        The \path|$MFT|, the change journal and \path|$LogFile| (the
        transaction journal, not converted) are also copied raw into the
        exhibit store, to rebuild the timeline later with another tool.
  \item[\fichier{recentdocs.json}] recent-document shortcuts: target, timestamps
        of the shortcut and of the target, target size, volume (type, serial
        number, label), network path, shell item list, and the extra data
        blocks: the tracker block (\champ{Tracker}: NetBIOS name of the machine
        the target was on, tracking identifiers, and what their version-1 GUID
        reveals — when the tracking began and the machine's hardware address),
        paths with environment variables, known folder, target properties.
  \item[\fichier{jumplistAutomaticDestinations.json},
        \fichier{jumplistCustomDestinations.json}] per-application jump lists,
        with the shortcuts they hold; in the automatic ones, each shortcut comes
        with its DestList entry (last access, machine name, tracking
        identifiers with their date and hardware address, pin status, number of
        openings \champ{InteractionCount}, rank \champ{MruPosition} — 0 for the
        most recent —, name of the known folder targeted).
  \item[\fichier{shellbags.json}] folders browsed in Explorer, as a rebuilt
        tree; for each folder, its rank among its siblings
        (\champ{MruPosition}, 0 for the most recent), the write time of the key
        holding it (\champ{LastWriteTime}) and of its own key
        (\champ{KeyLastWriteTime}: its addition, or the browsing of a
        subfolder), and its \champ{NodeSlot}. The ``first'' and ``last
        interaction'' of other tools are only readings of these two times.
  \item[\fichier{mrus.json}, \fichier{mruApps.json}] most-recently-used files
        and applications.
  \item[\fichier{zone\_identifiers.json}] \fichier{Zone.Identifier} streams
        (``Mark of the Web'') that browsers and mail clients write on what they
        download: file, zone (\champ{ZoneId}, \champ{ZoneName} —
        \texttt{Internet} for~3), page of origin (\champ{ReferrerUrl}) and the
        file's address (\champ{HostUrl}), every value in \champ{Values}. Read in
        the copy of the \path|$MFT|, where this small stream is resident: no
        further read of the examined machine; a non-resident stream, rare, is
        flagged (\champ{ContentNotResident}).
  \item[\fichier{recycle\_bin.json}] files sent to the recycle bin, on every
        fixed NTFS volume, from their \path|$I| index: original path, size,
        instant of the deletion (\champ{DeletedUtc}), and the user
        (\champ{SID}). Only the \path|$I| are collected; the files themselves
        (\path|$R|) are in \fichier{mft.json}.
\end{description}

\subsection{Devices}

\begin{description}[style=nextline,leftmargin=1.2em]
  \item[\fichier{setupapi.json}] sections of \fichier{setupapi.dev.log} and its
        rotations: operation (\champ{Operation}), device or driver
        (\champ{Subject}, with a USB key's serial number), start and end
        converted from the examined machine's local time, outcome. The first
        connection of a device is dated there, even when the \fichier{USBSTOR}
        key was cleaned.
  \item[\fichier{Usbstor.json}] USB storage devices known to the system:
        \champ{Vendor}, \champ{Product}, \champ{Revision} and \champ{Type},
        from \champ{DeviceId}; \champ{SerialNumber}, from \champ{InstanceId}
        — absent when Windows made the identifier up (second character
        ``\&''), the device declaring none; first and last insertion dates.
  \item[\fichier{mounted\_device.json}] mapping between drive letters and
        devices. \champ{Drive} names the letter or the volume; \champ{Type}
        gives the form of the value: \texttt{GPT partition}
        (\champ{PartitionGuid}), \texttt{MBR partition}
        (\champ{DiskSignature} and \champ{PartitionOffset}, in bytes),
        \texttt{Device path} (\champ{Device}, the device path, with the serial
        number of a USB key) or \texttt{Unknown} (\champ{Data}, the bytes in
        hexadecimal).
\end{description}

\subsection{Event logs (\opt{events})}

\fichier{events.json} holds one object per event:

\begin{description}[style=nextline,leftmargin=1.2em]
  \item[\champ{EvtSystem*}] the system section: provider, identifier, level,
        task, keywords, creation time, record number, process and thread,
        channel, computer, user;
  \item[\champ{EvtEventData}] the event's data, as a list of \champ{Name} /
        \champ{Value} pairs;
  \item[\champ{EvtEventMessage}] the readable message, as Event Viewer would show
        it, rebuilt from the provider's resources collected on the machine;
  \item[\champ{EvtSourceLog}] the log file the event comes from. One channel can
        be carried by several files (current log and archives) whose record
        numbers overlap;
  \item[\champ{EvtRecovered}] \texttt{true} for an event of a \emph{cleared}
        log: Windows resets the chunk header without wiping the records, which
        WAC reads again. Event Viewer no longer shows them; they remain the
        trace of what the clearing meant to erase, and must never be read as
        live events.
\end{description}

\begin{note}[Incomplete messages]
A \texttt{\%1}, \texttt{\%2}\ldots{} mark left in a message means the event does
not carry the expected data: it is kept visible so that the sentence does not
look complete. Not every provider ships message resources: their events then
have no \champ{EvtEventMessage}, but their data are complete.
\end{note}

% ============================================================ 6
\chapter{Reading the data}\label{chap:lecture}

\section{Timestamps}

Every date is emitted twice: \champ{X} in local time and \champ{XUtc} in UTC.
The local time is the \textbf{suspect's}, computed with the time zone read from
its \fichier{SYSTEM} hive — not the collecting machine's. Dates carry the
fraction of a second when the source provides it, and only to that source's
precision: seven digits for a \texttt{FILETIME}, three for a date to the
millisecond, none for a date to the second (FAT date, precise to two seconds,
text date, Unix time).

Every date carries the offset in force \textbf{at that date}, daylight saving
time included: a January date collected in summer in Paris is at
\texttt{+01:00}, a July date at \texttt{+02:00}. \wac{} applies the rules of
the examined machine year by year (\fichier{Dynamic DST} key of the
\fichier{SOFTWARE} hive): a 2005 date in the United States follows the rules
of 2005. The local time written is always derived from the UTC time: one instant
reads the same in every artefact. Two dates have no knowable time zone and
are given \textbf{as written}, without offset nor UTC, with \champ{TimeZone}
saying so: the date of an item seen inside an archive
(\champ{ModifiedAsWritten}; a ZIP entry keeps the clock time of the machine
that made the archive) and an OLE date of a property store
(\champ{AsWritten}). Placing them on a timeline needs their origin: \wac{}
does not make it up. For a value the artefact stores in local
time (FAT volume, \fichier{setupapi.dev.log}), the hour repeated in autumn is read as its first occurrence, and the
hour skipped in spring with the offset before the change, as Windows does.

\section{Account names}

The fields that name an account (\champ{SIDName}, \champ{Owner},
\champ{RunAs}) come from the evidence alone, without asking the system or
the network: local accounts and groups of the SAM database (in the language
of the examined machine), well-known SIDs under their canonical names
(\texttt{SYSTEM}, \texttt{LOCAL SERVICE}\ldots), service SIDs
(\texttt{NT SERVICE\textbackslash\ldots}), and the profile folder for a
domain account. A SID none of them names is published without a name: the SID
remains the evidence, for the analyst to resolve against the domain.

\section{GUID names}

The fields that name a GUID (COM class, known folder, handler of a task) come
first from a reference table, then from the examined machine's hives: its
classes (\fichier{SOFTWARE\textbackslash Classes\textbackslash CLSID}), its
known folders (\fichier{FolderDescriptions}), and each user's classes
(\fichier{UsrClass.dat}). A component installed on the machine is therefore
named, with the module that serves it. A GUID nothing names is published
without a name.

\section{A missing field is not an empty field}

A value with no source is not emitted at all. An empty key would read as a
failed read, and could not be told apart from ``the hive does not hold this
value''. What could not be read is recorded in \fichier{investigation.json}.

\section{Artefact not collected}

When a whole artefact could not be read, its file says so explicitly (the text
is produced in French):

\begin{lstlisting}
{
  "Artifact": "USBSTOR (registre)",
  "CollectionStatus": "NotCollected",
  "Error": "0x02 Le fichier spécifié est introuvable.",
  "Note": "La lecture de cet artefact a échoué : l'absence de données
           ci-dessus ne signifie PAS qu'aucune trace n'existe sur le système."
}
\end{lstlisting}

\begin{attention}
An empty artefact file and a ``NotCollected'' file do not say the same thing.
The first means no entry was found; the second, that the source could not be
read.
\end{attention}

\section{Undecoded data}

An item \wac{} cannot decode (unknown shell item, unrecognised extension block)
comes with its raw bytes in hexadecimal, so that an analyst can decode it later.
With \opt{dump}, the raw value of shellbags and shortcuts is always added.

\section{Paths}

Paths and values are delivered \textbf{as they are on the disk}, untransformed:
escaping happens only when the JSON is written.

% ============================================================ 7
\chapter{Traces left on the examined machine}\label{chap:traces}
\markboth{Traces on the examined machine}{}

Read this before using \wac{} on a case. This chapter lists what \wac{} does to
the machine, operation by operation, including what it cannot avoid.

\section{What WAC still does}

\begin{longtable}{@{}L{4.3cm}L{10.1cm}@{}}
\toprule
\bfseries Operation & \bfseries Trace \\
\midrule
\endhead
Direct volume reading & One access per volume actually read (usually one). No
  file opened, so no access date modified. If object-access auditing is enabled,
  opening the volume may be logged (Security 4656/4663). \\
Writing the results & On the collection medium only. Nothing is written to the
  examined disk. \\
Replaying hive logs & On the working copy only, with an undo journal. \\
Service state & One read-only query to the Service Control Manager. No service
  is opened or modified. \\
Processes & A kernel snapshot; no process is opened. Owners: one query to the
  Terminal Services service. \\
Sessions & One read-only query to LSASS. \\
Network state & Reads of the TCP/IP stack's tables in memory: no packet sent,
  no name resolved. \texttt{iphlpapi.dll} is loaded only when
  \texttt{artefacts.network} is on. \\
Boot time & One in-memory query to the kernel. \\
Event logs (\opt{events}) & Read through the same direct volume access: the
  logging service is not solicited. \\
Referenced files (\opt{binary}) & Read through the same direct access. No file
  opened. \\
Authenticity check (\opt{binary}) & Catalogs read through the same direct
  access; verification in memory. No service solicited, no access to the
  certificate store, no network access. \\
\bottomrule
\end{longtable}

\section{What WAC no longer does}

Earlier versions created a Volume Shadow Copy, went through COM and WMI, queried
the Task Scheduler and the logging service, opened every process and every
service, read the profile list in the live registry and opened every referenced
file to fingerprint it. All these operations, which left entries in the logs or
changed access dates, have been removed.

\section{What no tool can avoid}

These traces are listed so that the analyst recognises \wac{}'s footprint
instead of taking it for the activity under study:

\begin{itemize}
  \item \textbf{running \fichier{WAC.exe} is an execution}: it produces a
        Prefetch file for \fichier{WAC.exe}, an Amcache entry and, if
        process-creation auditing is on, a Security 4688 event. Expect to find
        \wac{} in the very artefacts it collects;
  \item \textbf{plugging in the USB stick} registers it in \fichier{USBSTOR},
        \fichier{MountedDevices} and the related logs;
  \item \textbf{logging on} to run the tool creates a session and its log
        entries;
  \item \wac{} \textbf{takes up memory} during the collection, which may push
        pages out to the page file. Its reading of the volume bypasses the
        cache (unbuffered reads): the gigabytes read do not replace the
        machine's file cache.
\end{itemize}

% ============================================================ 8
\chapter{Known limitations}

\begin{description}[style=nextline,leftmargin=1.2em]
  \item[Volumes neither NTFS nor FAT] A file on a CD-ROM is not read: it is
        reported as unread.
  \item[Per-user variables] A path written with \path|%APPDATA%|,
        \path|%LOCALAPPDATA%| or \path|%USERPROFILE%| (in a scheduled task, for
        instance) is not resolved: the account concerned is not known there, and
        guessing would fingerprint another file.
  \item[Deleted executable] A process whose executable has been deleted from the
        disk has no fingerprint: only memory is left, outside \wac{}'s scope.
  \item[Event messages] About a third of the events get a readable message on
        the test machine; the other providers ship no usable resources. The data
        themselves are always complete.
  \item[Signed but not authenticated] A few kinds of genuinely signed files are
        still collected, the cautious side: Microsoft Store applications (signed
        at package level, not file by file) and Windows Script Host scripts
        signed only within their text.
\end{description}

% ============================================================ 9
\chapter{Troubleshooting}

\begin{longtable}{@{}L{5cm}L{9.4cm}@{}}
\toprule
\bfseries Message or finding & \bfseries Cause and remedy \\
\midrule
\endhead
``Check administrator rights'' fails & The prompt was not started as
  administrator. Start it again with administrator rights. \\
``Check OS >= Windows 10'' fails & System older than Windows~10: not
  supported. \\
``Checking the collection medium'' fails, directory not empty & The output
  folder holds an earlier collection. Give a new folder with \opt{output}. \\
``Checking the collection medium'' fails, not enough space & The medium is too
  small for the chosen options (section~\ref{sec:place}). Use a larger medium or
  drop \opt{binary}. \\
An artefact file holds \texttt{"CollectionStatus": "NotCollected"} & The source
  could not be read; the reason is in \champ{Error}. The other artefacts are not
  affected. \\
Failed exhibits in the manifest & Referenced files missing from the disk
  (uninstalled software, service profile without a hive), non-NTFS
  volume\ldots{} The reason is in \champ{Result}. \\
The log reports binaries ``hashed without being copied'' & Space ran out during
  collection: the fingerprints are complete, the copies are not. \\
Need to understand a read failure & Run again with \opt{loglevel=2} (and, for
  NTFS reading, \opt{debug}) and read \fichier{WAC.exe.log}. \\
\bottomrule
\end{longtable}

% ============================================================ APPENDICES
\appendix

\chapter{Glossary}

\begin{description}[style=nextline,leftmargin=1.2em]
  \item[Artefact] Trace left by the activity of the system or of a user, usable
        in an investigation.
  \item[Exhibit store (\fichier{exhibits})] Directory of the raw copies, sealed,
        never rewritten.
  \item[Fingerprint] Value computed on a file's content (MD5, SHA-1, SHA-256):
        different contents have different fingerprints.
  \item[Hive] File holding part of the registry (\fichier{SYSTEM},
        \fichier{SOFTWARE}, \fichier{ntuser.dat}\ldots).
  \item[Manifest] File identifying each exhibit of the exhibit store.
  \item[\path|$MFT|] File table of an NTFS volume.
  \item[Seal] The manifest's fingerprint, kept apart: it reveals any change to
        the manifest.
  \item[Transaction log] \fichier{.LOG1}/\fichier{.LOG2} file holding a hive's
        recent changes, not yet written into it.
  \item[Working copy (\fichier{working})] Directory of the working copies, on
        which the analysis and any modification are made.
  \item[WOF / Compact OS] Compressed storage mode of the system files of
        Windows~10 and~11.
\end{description}

\end{document}
